\section{Data}
\label{sec:data}

Since our central claim is that generalization lives in the training data
rather than the architecture, the data deserves the same scrutiny usually
reserved for the model.

\subsection{Training mixture}

We assemble 279 tasks and 323{,}466 rows from
\texttt{tasksource}~\cite{tasksource}, normalised into the typed-decision
format. Composition after cleaning:

\begin{center}
\begin{tabular}{lrl}
\toprule
primitive & tasks & notes \\
\midrule
\choice{} & 234 & option counts 2 to 174; 80 tasks have per-example menus \\
\noul{}   & 34  & \\
\score{}  & 11  & mostly three-level sentiment; see \S\ref{sec:discussion} \\
\midrule
total & 279 & 323{,}466 rows \\
\bottomrule
\end{tabular}
\end{center}

The lopsidedness is a limitation, not a design choice, and it predicts the
result pattern: \choice{} transfers best and \score{} worst.

\paragraph{Per-example option menus.} 80 tasks come from the
multiple-choice family, where each row carries its own answer set rather
than sharing a fixed menu. Supporting this in the trainer is what unlocked
that family, which is 261 of the 668 tasks \texttt{tasksource} exposes.

\subsection{Contamination guard}

A held-out suite is worthless if the training mixture contains the same
data under a different name, and aliasing is the normal case:
\texttt{tasksource} calls one dataset \texttt{banking77}, while the Hub
also carries \texttt{PolyAI/banking77} and \texttt{mteb/banking77}.

Each evaluation task therefore declares every alias of its source, a
manifest holds the union, and a matcher checks full names, owner-stripped
basenames, and glob patterns rather than performing a set intersection.
The check runs before every training run and refuses to proceed on a hit;
it excluded 23 tasks when building this mixture, and caught a real leak
during development. We verify it with a positive control --- injecting
\texttt{banking77} into a clean mixture and confirming rejection --- since
a guard that has never fired is indistinguishable from a guard that
cannot.

\subsection{Auditing, and what it found}
\label{sec:data-audit}

Training is the expensive step, so faults that survive it and surface
hours later as a bad number are worth catching first. We audit for gold
answers outside their own menu, \score{} golds stored as booleans (in
Python \texttt{True == 1}, which silently rewrites levels 0 and 1), class
imbalance too severe to learn, near-duplicate \noul{} descriptions, split
leakage, and contradictory labels.

Three findings changed the data materially.

\paragraph{The gold answer was always first.} \texttt{tasksource} ships
multiple-choice tasks correct-answer-first. Measured after ingestion, the
gold sat at index $0$ in \textbf{100\% of rows across all 83 ingested
tasks}. Stored that way, the data teaches ``pick the first option'' to
anything reading the menu as written. We shuffle at ingestion so the
stored data is honest, and drop tasks whose gold remains at a fixed index
or is a single constant string.

\paragraph{Contradictory labels.} 1{,}994 rows carried the same state under
the same menu with different answers; whichever the model predicts, some
are wrong, so the gradient is noise. One task contributed 195.

\paragraph{Menus too long to pack.} 1{,}787 rows had option sets exceeding
the context budget on their own, some options being multi-paragraph
passages.

Finally, we sweep every example through the same dataset object the
trainer uses, including augmentation, before training starts: 323{,}466
examples verified in 227 seconds of CPU. Two earlier runs died on packing
errors after twenty minutes and ten thousand steps respectively, and both
faults were findable in advance.

\subsection{Label-space augmentation}
\label{sec:data-aug}

Our mixture's option counts topped out near 20 while the evaluation suite
reaches 151. A model that has never seen a large menu collapses onto one
label when handed one, which is what an exact $0.000$ looks like: chance
alone would have produced roughly three correct in 200.

The intervention is to manufacture the missing shape. With probability
$\pi$, a training menu is padded with labels borrowed from other tasks in
the mixture, sampling a target size log-uniformly so small menus stay
common and large ones appear often enough to matter. The gold answer is
unchanged and still correct, so the example remains valid, but the model
must now discriminate against a large menu.

\begin{algorithm}[t]
\caption{Label-space augmentation for one \choice{} question}
\label{alg:aug}
\begin{algorithmic}[1]
\Require menu $O$ with gold $y$, distractor pool $D$, probability $\pi$,
  cap $M$
\If{$\textproc{Uniform}(0,1) \ge \pi$} \Return $O$ \EndIf
\State $n \gets \big\lfloor \exp\big(\textproc{Uniform}(\ln(|O|{+}1),\,
  \ln M)\big) \big\rfloor$ \Comment{log-uniform target size}
\State $E \gets \emptyset$
\While{$|O| + |E| < n$ and attempts remain}
  \State $d \sim D$
  \If{$d \notin O$ and $d \notin E$} \Comment{never add a plausibly correct label}
    \State $E \gets E \cup \{d\}$
  \EndIf
\EndWhile
\State \Return $\textproc{Shuffle}(O \cup E)$ \Comment{$y$ still present and correct}
\end{algorithmic}
\end{algorithm}

Option order is shuffled every epoch for a second reason: on the reference
system we measured the argmax flipping across option orderings on 7.5\% of
confusable items, at 2.8$\times$ the resampling floor. That is per-item
instability rather than a positional prior, so it cannot be corrected post
hoc and must be trained out.

\subsection{Evaluation suites}

\paragraph{Held-out suite (Study G).} Seven public tasks, 600 rows each,
never trained on: \texttt{clinc\_oos} ($K{=}151$), \texttt{banking77}
($K{=}77$), \texttt{massive\_intent} ($K{=}60$), \texttt{sst5} and
\texttt{helpsteer} (\score{}, $K{=}5$), \texttt{ag\_news} ($K{=}4$), and
\texttt{civil\_comments} (\noul{}). The high-cardinality end is the point:
a 151-way menu is where models that look fine on four-way classification
fail.

Every evaluation reports a \emph{held-in control}: accuracy on tasks the
model was trained on. If held-out is at chance and held-in is high, the
model genuinely fails to transfer; if both are at chance, it is a bug. We
also run a trivial three-way sanity task that any working harness must
pass. This is not ceremony --- the sanity check caught a real
partially-loaded-model bug that had produced a plausible, and wrong,
results table.

\paragraph{Healthcare router (Study S).} A synthetic pharmacy routing task:
988 items across 18 tiers, 395 train and 450 test, with ten questions per
item (one intent \choice{}, five topic \noul{}s, four safety gates).

Its first version was useless: the reference API scored a clean $1.000$ on
four of its tiers, which cannot distinguish any two systems. We rebuilt it
targeting where that system actually failed --- obliquely-worded clinical
risk, three-intent compounds, negation and conditionals, near-miss
out-of-scope, subtle prompt injection, and degraded voice-transcript text.
Splits are disjoint by content rather than by row, so paraphrases and
noisy variants of an item all stay in one split.
